For  x^n + a_{n-1}x^{n-1} + \cdot\cdot\cdot + a_1 x + a_0 = 0  with roots  r_1, r_2, \cdot\cdot\cdot, r_n:

\sum r_i = -a_{n-1}
\quad\quad (sum of roots = -coeff of  x^{n-1})

\sum_{i < j} r_i r_j = a_{n-2}
\quad\quad (sum of products of pairs)

r_1 r_2 \cdot\cdot\cdot r_n = (-1)^n a_0
\quad\quad (product of all roots)

For quadratic  ax^2 + bx + c = 0:
r_1 + r_2 = -\frac{b}{a}, \quad r_1 r_2 = \frac{c}{a}

For cubic  ax^3 + bx^2 + cx + d = 0:
r_1+r_2+r_3 = -\frac{b}{a}
r_1 r_2 + r_1 r_3 + r_2 r_3 = \frac{c}{a}
r_1 r_2 r_3 = -\frac{d}{a}

Sum of squares of roots:
r_1^2+r_2^2+r_3^2 = (r_1+r_2+r_3)^2 - 2(r_1r_2+r_1r_3+r_2r_3)

Absolute difference of roots (quadratic):
|r_1 - r_2| = \frac{\sqrt{b^2 - 4ac}}{|a|}

Reciprocal roots: if  r  is a root of  f(x), then  \frac{1}{r}
is a root of  x^n f\!\left(\frac{1}{x}\right)
\rightarrow reverse the coefficient list

Conjugate pairs: complex roots of real-coefficient polynomials
come in conjugate pairs.

Common root of two polynomials: subtract to reduce degree,
then find the root.

Example (2024 \#8):  x^3+kx^2-3x+4=0  and  x^3+kx^2-5x+8=0
Subtract:  2x - 4 = 0 \rightarrow x = 2
Sub back:  8 + 4k - 10 + 8 = 0 \rightarrow k = -\frac{3}{2}

Remainder Theorem:  f(x) \div (x - c)  has remainder  f(c).

Factor Theorem:  (x-c)  is a factor \rightarrow  f(c) = 0

Rational Root Theorem: possible rational roots = \pm\frac{factors of constant}{factors of leading coeff}

Polynomial with given zeros  z_1, z_2, \cdot\cdot\cdot, z_n  (leading coeff 1):
f(x) = (x - z_1)(x - z_2) \cdot\cdot\cdot (x - z_n)
